\documentclass[10pt,a4paper]{amsart}
\usepackage[margin=19mm]{geometry}
\usepackage{amsmath,amssymb,mathtools}
\usepackage[scr=rsfs,cal=euler]{mathalfa}
\usepackage{xfrac}
\usepackage[unicode,hidelinks,bookmarks=false]{hyperref}
\newcommand{\ie}{i.e.\ }
\newcommand{\cf}{cf.\ }
\newcommand{\resp}{respectively\ }
\newcommand{\Subsection}{\subsection}
\providecommand{\nobreakdash}{\nobreak\hbox{-}}
\setlength{\emergencystretch}{2em}
\multlinegap0pt
\usepackage{amssymb}
\usepackage{xcolor}
\usepackage{algorithm}\usepackage{algpseudocode}
\usepackage{bm}
\usepackage{float}
\usepackage{tikz}
\newtheorem{theorem}{Theorem}
\newtheorem{corollary}{Corollary}
\newcommand{\E}{\mathbb{E}}
\newcommand{\Var}{\mathrm{Var}}
\newcommand{\bea}{\begin{eqnarray}}
\newcommand{\ena}{\end{eqnarray}}
\title{An approximate zero bias transformation for random sums: Applications to sampling with outliers, auto insurance, and generative AI}
\author{Wasamon Jantai, Nathakhun Wiroonsri}
\date{Source: arXiv:2608.27143v1 (CC BY 4.0)}
\begin{document}
\maketitle

\noindent\emph{Source and licence.} An approximate zero bias transformation for random sums: Applications to sampling with outliers, auto insurance, and generative AI. Version arXiv:2608.27143v1 (CC BY 4.0), distributed by arXiv under the Creative Commons Attribution 4.0 International licence, http://creativecommons.org/licenses/by/4.0/, as recorded in the arXiv licence metadata for this exact version and shown on the abstract page https://arxiv.org/abs/2608.27143v1. Full licence notice: \texttt{licenses/arxiv-2608.27143-CC-BY-4.0.txt}.\par\smallskip

\noindent\textit{Editorial scope.} Counted unit: the article's corollary Cor-outliers1 (source0.tex lines 680--683). The article's complete original proof of it is delivered with it (source0.tex lines 684--699, one \texttt{proof} environment of 16 lines). No mathematics is written, repaired or re-derived: the statement and the proof are the article's own lines, in the article's own notation, and no retained line is substituted or re-flowed.  Retained uncounted as the source-local context the proof needs: lines 306--325; lines 462--469.  Nothing was omitted from any retained span, so the package carries no omission record.  No retained line contains a citation, so the wrapper carries no bibliography and no bibliography entry.  No retained line contains a figure environment, an image-inclusion command or drawing code, so no image is delivered and the package declares no asset dependency.  Editorial formatting only, in addition: an AMS \texttt{amsart} preamble that restates the article's own declarations and macros used by the retained lines, namely the environments \texttt{theorem}, \texttt{corollary} on separate counters, together with the standard packages those lines need.  The retained environments print in this document as Theorem 1, Corollary 1 in order of appearance, whereas the article prints the same items with the numbering of its own full document.  The complete native source of the article as distributed in the arXiv e-print for this exact version is bundled unchanged as \texttt{sources/208691/source0.tex}.
\textbf{
	\begin{flushleft}
		Construction of an approximate zero bias variable:
	\end{flushleft}
} \label{construction}
\begin{enumerate}
	\item For $i\in [n]$, let $X_i^*$ has the zero bias distribution of $X_i$ independent of $(X_j)_{j\ne i}$ and $(X_j^*)_{j\ne i}$.
	\item Choose a random summand $X_I$ with 
	\bea \label{Idef}
	P(I=i) = \frac{\sigma_i^2P(N'\ge i)}{\sigma^2},
	\ena  
	where $N'$ has the same distribution as $N$ and independent of $N$ and $I$ is independent of all other random variables.
	\item Define 
	\bea \label{Wsdef}
	W^* = \begin{cases}
		\sum_{j\ne I}X_j+X_I^*, &\text{ \ if \ } N \ge I,\\
		W,  &\text{ \ if \ } N<I.
	\end{cases}
	\ena
\end{enumerate}

\begin{theorem}\label{uniform}
	Let $X_1,X_2,\ldots$ be independent mean zero random variables such that $\Var(X_i) = \sigma_i^2$. Let $W=\sum_{i=1}^NX_i$, where $N$ is a discrete random variable with support being a subset of $\mathbb{N}_0$ and independent of all $X_i$.
		Define $\sigma^2 := \Var(W) = \sum_{i=1}^\infty\sigma_i^2 P(N\ge i)$. Let $I$ be as in \eqref{Idef}.  Then
	\begin{align*}
		\left|\E h\left(\frac{W}{\sigma}\right)-N(h)\right|\leq \frac{||f'_h||}{\sigma^3}\sum_{i=1}^\infty \left({1\over 2}\E|X_i|^3+ \sigma_i^2\E|X_i| \right)P^2(N\ge i)+2 ||f^{'}_{h}|| \cdot {P(N < I)},
	\end{align*}
	where $N(h) = \E h(Z)$ with $Z \sim \mathcal{N}(0,1)$, $h \in \mathcal{H}_1$ and $f_h$ is the solution to the Stein's equation.
	\end{theorem}

\begin{corollary}\label{Cor-outliers1}
	Let $W$ be a mean zero, variance $\sigma^2 >0$ random variable, then
	\begin{align*} &d_1\left(\mathcal{L}(W/\sigma), \mathcal{L}(Z) \right) \\&\quad \le\left({\left( \dfrac{12}{e}-2\right) \over \left({nN\over N+m}\right)^{3/2}} +{\left( \dfrac{4}{e}\right)\over \left({nN\over N+m}\right)^{3/2}}\right)\sum_{i=\max\{0,n-m\}}^{\min\{n,N\}} P^2(M\ge i)  \\&\quad \quad+{2\sqrt{2} \over \sqrt{\pi}}\cdot\displaystyle{\left(1- \left( \frac{N+m}{nN}\right)\sum_{i=\max\{0,n-m\}}^{\min\{n,N\}} P^2(M\ge i)\right)},\end{align*} where $Z$ is a standard normal random variable.
\end{corollary}

\begin{proof}
	Using Theorem~\ref{uniform} together with the bounds 
	$\|f_h''\| \le 2$ and $\|f_h'\| \le \sqrt{2/\pi}$, and the identity
	\begin{align*}
		P(M \ge I) 
		&= \sum_{i=\max\{0,n-m\}}^{\min\{n,N\}}  \frac{\sigma_i^2}{\sigma^2} \, P(M \ge i)^2\\&=  \sum_{i=\max\{0,n-m\}}^{\min\{n,N\}} \frac{r^2}{\sigma^2} \, P(M \ge i)^2\\& =  \left( \frac{N+m}{nN}\right)\sum_{i=\max\{0,n-m\}}^{\min\{n,N\}} \, P(M \ge i)^2,
	\end{align*} 
	we obtain \begin{align*} &d_1\left(\mathcal{L}(W/\sigma), \mathcal{L}(Z) \right) \\&\quad \le \left({\E|X_1|^3 \over \left({nNr^2\over N+m}\right)^{3/2}} +{2r^2\E|X_1|\over \left({nNr^2\over N+m}\right)^{3/2}}\right)\sum_{i=\max\{0,n-m\}}^{\min\{n,N\}} P^2(M\ge i)  \\&\quad \quad+{2\sqrt{2} \over \sqrt{\pi}}\cdot\displaystyle{\left(1- \left( \frac{N+m}{nN}\right)\sum_{i=\max\{0,n-m\}}^{\min\{n,N\}} P^2(M\ge i)\right)}.\end{align*}
	For $X\sim \mathrm{Exp}(\lambda)$, note that $\lambda X\sim \mathrm{Exp}(1)$. Using the identities
	\[
	\E|X-\E X|^3 = \frac{1}{\lambda^3}\left(\frac{12}{e}-2\right), \qquad
	\E|X-\E X|^2 = \frac{1}{\lambda^2}, \qquad
	\E|X-\E X| = \frac{2}{\lambda e},
	\]
	the desired expression follows.
\end{proof}

\end{document}
